\documentclass[12pt]{article}
\usepackage{geometry}
\geometry{margin=1in}

% ----------------------------------------------------------
%  Linux Penguin Theme — Based on Recommended Colors
% ----------------------------------------------------------

% --- COLOR SETUP ---
\usepackage{xcolor}

\definecolor{cream}{RGB}{246,242,219}
\definecolor{teamred}{RGB}{209,57,64}
\definecolor{navy}{RGB}{40,46,87}
\definecolor{softwhite}{RGB}{250,250,230}
\definecolor{accentyellow}{RGB}{240,200,60}

% Extra utility colors
\definecolor{muted}{RGB}{120,125,140}
\definecolor{darkcream}{RGB}{230,225,205}

% ----------------------------------------------------------
%  FONT SETUP (Optional but Recommended)
%  Works best with XeLaTeX or LuaLaTeX
% ----------------------------------------------------------
\usepackage{fontspec}

% Main document font (Google fonts — upload to Overleaf)
% If you don't upload fonts, comment these two lines.
\setmainfont{SourceSans3}[
  Path=./,
  ItalicFont  = *-Italic,
  BoldFont    = *-Bold
]

\newfontfamily\headingfont{ScienceGothic}[
  Path=./,
]

% Font Awesome for minimal icons
\usepackage{fontawesome5}

% ----------------------------------------------------------
%  SECTION STYLE
% ----------------------------------------------------------
\usepackage{titlesec}

% Section
\titleformat{\section}
  {\Large\headingfont\bfseries\color{navy}}
  {\thesection}
  {1em}
  {\rule{\linewidth}{1pt}\vspace{0.5em}\\}

% Subsection
\titleformat{\subsection}
  {\large\headingfont\bfseries\color{teamred}}
  {\thesubsection}
  {0.8em}
  {}

% Subsubsection
\titleformat{\subsubsection}
  {\bfseries\color{navy}}
  {\thesubsubsection}
  {0.6em}
  {}

% ----------------------------------------------------------
%  PAGE STYLE
% ----------------------------------------------------------
\usepackage{fancyhdr}
\pagestyle{fancy}

\fancyhf{} % clear all

\lhead{\textcolor{navy}{\headingfont \courseTitle}}
\rhead{\textcolor{teamred}{\thepage}}
\renewcommand{\headrule}{\color{teamred}\hrule height 1pt}

% ----------------------------------------------------------
%  HYPERLINK STYLE
% ----------------------------------------------------------
\usepackage{hyperref}

\hypersetup{
  colorlinks = true,
  linkcolor  = teamred,
  urlcolor   = navy,
  citecolor  = accentyellow
}

% ----------------------------------------------------------
%  CUSTOM BOX ENVIRONMENTS
% ----------------------------------------------------------
\usepackage{tcolorbox}
\tcbuselibrary{skins,breakable}

% Info box
\newtcolorbox{infobox}{
  enhanced,
  breakable,
  colback   = softwhite,
  colframe  = navy,
  coltext   = navy,
  boxrule   = 1pt,
  arc       = 4pt,
  left=6pt, right=6pt, top=6pt, bottom=6pt
}

% Warning box
\newtcolorbox{warningbox}{
  enhanced,
  breakable,
  colback   = accentyellow!30,
  colframe  = accentyellow!80!navy,
  coltext   = navy,
  boxrule   = 1pt,
  arc       = 4pt,
  left=6pt, right=6pt, top=6pt, bottom=6pt
}

% Highlight box
\newtcolorbox{highlightbox}{
  enhanced,
  breakable,
  colback   = teamred!10,
  colframe  = teamred,
  coltext   = navy,
  boxrule   = 1pt,
  arc       = 4pt,
  left=6pt, right=6pt, top=6pt, bottom=6pt
}

% ----------------------------------------------------------
%  ITEMIZE / ENUMERATE STYLE
% ----------------------------------------------------------
\usepackage{enumitem}

\setlist[itemize]{
  label=\textcolor{teamred}{\large\textbullet},
  leftmargin=1.2em,
  itemsep=4pt
}

\setlist[enumerate]{
  label=\textcolor{navy}{\arabic*.},
  leftmargin=1.4em,
  itemsep=4pt
}

% ----------------------------------------------------------
%  TITLE STYLE
% ----------------------------------------------------------
\usepackage{titling}

\pretitle{
  \begin{center}
  \headingfont\bfseries\LARGE\color{navy}
}
\posttitle{
  \end{center}
}

\preauthor{
  \begin{center}
  \color{teamred}
}
\postauthor{
  \end{center}
}

% ----------------------------------------------------------
%  OPTIONAL — SOFT PAGE BACKGROUND
%  Comment out if you want white pages
% ----------------------------------------------------------
\usepackage{pagecolor}
\pagecolor{cream}
\color{navy}

% ----------------------------------------------------------
% CUSTOM COMMANDS FOR TEMPLATE
% ----------------------------------------------------------
\newcommand{\courseTitle}{Linux Process Monitoring \& Management}
\newcommand{\courseDate}{23-02-2026}
\newcommand{\authorName}{Khaled Mahfouz}
\newcommand{\contactInfo}{\texttt{@khaledmhfz}}

% Minimal icon commands
\newcommand{\penguinicon}{\faLinux}
\newcommand{\htopicon}{\faChartLine}
\newcommand{\killicon}{\faSkullCrossbones}
\newcommand{\searchicon}{\faSearch}
\newcommand{\backgroundicon}{\faSleep}
\newcommand{\chainicon}{\faLink}
\newcommand{\summaryicon}{\faList}
\newcommand{\tipsicon}{\faLightbulb}
\newcommand{\bookicon}{\faBook}
\newcommand{\arrowicon}{\faArrowRight}

% ----------------------------------------------------------
% END OF THEME
% ----------------------------------------------------------

\begin{document}

% Title with minimal space
\begin{center}
\headingfont\Huge\bfseries\color{navy}
\penguinicon\ \courseTitle\\[0.5em]
\Large\color{teamred}
\courseDate
\end{center}

\vspace{1em}
\begin{center}
\color{teamred}\rule{0.8\textwidth}{1pt}
\end{center}

\section*{\faInfoCircle\ What Is a Process?}

Every program you run becomes a \textit{process} — an instance of a running program with a unique ID (PID). Linux lets you \textbf{list}, \textbf{monitor}, and \textbf{control} these processes using simple tools.

\section*{\htopicon\ 1) \textbf{htop — Interactive Process Viewer}}

\begin{infobox}
\faLightbulb\ \texttt{htop} is like \texttt{top}, but \textit{way more human-friendly}: color, scrolling, arrow navigation, signal menus, and killer shortcuts. (\href{https://www.geeksforgeeks.org/linux-unix/using-htop-to-monitor-system-processes-on-linux/}{\texttt{GeeksforGeeks}})
\end{infobox}

\subsection*{Why \texttt{htop}?}

\begin{itemize}
    \item Shows CPU, RAM, threads in real-time
    \item Scroll through processes
    \item Select and kill without typing PIDs
\end{itemize}

\subsection*{Install \& Run (Debian/Ubuntu/Mint)}

\begin{highlightbox}
\texttt{\$ sudo apt update}\\
\texttt{\$ sudo apt install htop}\\
\texttt{\$ htop}
\end{highlightbox}

\subsection*{Inside \texttt{htop}}

\begin{itemize}
    \item \faArrowUp\ / \faArrowDown\ — move
    \item F9 — kill selected process
    \item F3 — search/filter
    \item F5 — tree view (parent/child structure)
\end{itemize}

\begin{infobox}
\arrowicon\ Best tutorial: \textit{GeeksforGeeks: Using htop to Monitor System Processes on Linux} (\href{https://www.geeksforgeeks.org/linux-unix/using-htop-to-monitor-system-processes-on-linux/}{\texttt{Click Here !!}})
\end{infobox}

\section*{\killicon\ 2) \texttt{kill} and \texttt{pkill} — End Processes}

When a process misbehaves (freezes, uses all your CPU), you kill it.

\subsection*{\texttt{kill}}

Targets a specific \textit{PID}.

\begin{highlightbox}
\texttt{\$ kill 1234}\\
\texttt{\$ kill -9 1234    \# force kill}
\end{highlightbox}

\begin{itemize}
    \item SIGTERM (\texttt{kill}) politely asks to quit
    \item SIGKILL (\texttt{kill -9}) stops immediately
\end{itemize}

\subsection*{\texttt{pkill}}

Targets by \textit{name} or pattern — no need to look up PID:

\begin{highlightbox}
\texttt{\$ pkill firefox}\\
\texttt{\$ pkill -u someuser}
\end{highlightbox}

\begin{infobox}
\faBrain\ \textbf{Tip:} \texttt{kill} works on PIDs, \texttt{pkill} works on \textit{names}.
\end{infobox}

\begin{infobox}
\arrowicon\ Great deep dive: \textit{How to Manage Linux Processes Using ps, kill, and pkill (How-To Geek)} (\href{https://www.howtogeek.com/how-to-manage-linux-processes-using-ps-kill-and-pkill/}{\texttt{Click Here !!}})
\end{infobox}

\section*{\searchicon\ 3) \texttt{ps} + \texttt{grep} — Find Processes via CLI}

\texttt{ps} gives you a \textit{snapshot} of currently running processes — it isn't real-time like \texttt{htop}, but it's extremely useful.

\subsection*{Common Pattern}

\begin{highlightbox}
\texttt{\$ ps aux            \# list ALL processes}\\
\texttt{\$ ps aux | grep firefox   \# find firefox}\\
\texttt{\$ ps -ef | grep python    \# find python instances}
\end{highlightbox}

\begin{infobox}
\faEye\  \texttt{grep} filters output — super handy when you're hunting a specific process. 
\end{infobox}

\begin{infobox}
\faBolt\ \textbf{Pro tip:} wrap your grep to \textit{not} match itself (e.g., \texttt{grep "[f]irefox"}), because normal grep often shows its \textit{own} process.
\end{infobox}

\begin{infobox}
\arrowicon\ For a beginner overview of \texttt{ps}, search, and manage processes, see \textit{Ubuntu Community Tutorial}. (\href{https://discourse.ubuntu.com/t/viewing-and-monitoring-processes-in-linux/26024}{\texttt{Click Here !!}})
\end{infobox}

\section*{4) Run in Background (\texttt{\&})}

Sometimes you want to start a long-running task \textit{without blocking your terminal}. That's where \textbf{background jobs} come in.

\subsection*{Example}

\begin{highlightbox}
\texttt{\$ sleep 60 \&}
\end{highlightbox}

\begin{itemize}
    \item \texttt{\&} → run command in the background
    \item You get your prompt back immediately
    \item Shell assigns a job number and PID
\end{itemize}

\subsection*{Job Management}

\begin{highlightbox}
\texttt{\$ jobs            \# list jobs}\\
\texttt{\$ fg \%1           \# bring job 1 to foreground}\\
\texttt{\$ bg \%1           \# resume job 1 in background}
\end{highlightbox}

\begin{infobox}
\faLocationArrow\ These are Bash's built-in \textit{job control} features.
\end{infobox}

\begin{infobox}
\arrowicon\ For a full walk-through on background/foreground jobs, see DigitalOcean's guide. (\href{https://www.digitalocean.com/community/tutorials/how-to-use-bash-s-job-control-to-manage-foreground-and-background-processes}{\texttt{Click Here !!}})
\end{infobox}

\section*{\chainicon\ 5) \textbf{Command Chaining — \texttt{\&\&}, \texttt{||}, \texttt{;}}}

Command chaining lets you fire off multiple commands in one line — \textit{with logic about success or failure}.

\subsection*{The Operators}

\begin{infobox}
\begin{tabular}{ll}
\textbf{Operator} & \textbf{What it Means} \\
\texttt{cmd1 ; cmd2} & Run cmd1, then cmd2 no matter what \\
\texttt{cmd1 \&\& cmd2} & Run cmd2 ONLY if cmd1 succeeds (exit status 0) \\
\texttt{cmd1 || cmd2} & Run cmd2 ONLY if cmd1 fails \\
\end{tabular}
\end{infobox}

\subsection*{Examples}

\begin{highlightbox}
\texttt{\$ sudo apt update \&\& sudo apt upgrade -y}\\
\texttt{\$ rm temp.log || echo "Could not delete temp.log"}\\
\texttt{\$ cmd1 ; cmd2 ; cmd3}
\end{highlightbox}

\begin{infobox}
\faCheckCircle\ Perfect for chaining updates, scripts, and error handling.
\end{infobox}

\begin{infobox}
\arrowicon\ See \textit{GeeksforGeeks: Chaining Commands in Linux} for examples. (\href{https://www.geeksforgeeks.org/linux-unix/chaining-commands-in-linux/}{\texttt{Click Here !!}})
\end{infobox}

\vspace{1em}
\begin{center}
\color{teamred}\rule{0.8\textwidth}{1pt}
\end{center}

\section*{\summaryicon\ Summary — Process Control Toolkit}

\begin{infobox}
\begin{tabular}{ll}
\textbf{Task} & \textbf{Tools Covered} \\
View all processes & \texttt{ps}, \texttt{htop}, real-time \\
Find specific process & \texttt{ps | grep} \\
Kill / stop process & \texttt{kill}, \texttt{pkill} \\
Run jobs & \texttt{\&}, \texttt{jobs} \\
Chain actions & \texttt{\&\&}, \texttt{||}, \texttt{;} \\
\end{tabular}
\end{infobox}

\section*{\bookicon\ Recommended Reading/Docs}

\begin{itemize}
    \item \penguinicon\ Ubuntu Command Line For Beginners — official Canonical tutorial (\href{https://ubuntu.com/tutorials/command-line-for-beginners}{\texttt{Ubuntu}})\\
    \arrowicon\ \href{https://ubuntu.com/tutorials/command-line-for-beginners}{\texttt{https://ubuntu.com/tutorials/command-line-for-beginners}}
    
    \item \penguinicon\ Bash Reference Manual — job control details\\
    \arrowicon\ \href{https://www.gnu.org/software/bash/manual/}{\texttt{https://www.gnu.org/software/bash/manual/}}
\end{itemize}

\subsection*{Reference Links}

\begin{itemize}
    \item \href{https://www.geeksforgeeks.org/linux-unix/using-htop-to-monitor-system-processes-on-linux/}{\texttt{Using htop to Monitor System Processes on Linux - GeeksforGeeks}}
    \item \href{https://www.howtogeek.com/how-to-manage-linux-processes-using-ps-kill-and-pkill/}{\texttt{How to Manage Linux Processes Using ps, kill, and pkill}}
    \item \href{https://www.owais.io/blog/2025-10-06_linux-process-management-monitoring-beginners/}{\texttt{Complete Guide to Linux Process Management and Monitoring for Absolute Beginners | Owais.io}}
    \item \href{https://www.reddit.com/r/linuxquestions/comments/dneenk}{\texttt{Grep command appearing when grepping processes?}}
    \item \href{https://discourse.ubuntu.com/t/viewing-and-monitoring-processes-in-linux/26024}{\texttt{Viewing and Monitoring Processes in Linux - Tutorials - Ubuntu Community Hub}}
    \item \href{https://www.digitalocean.com/community/tutorials/how-to-use-bash-s-job-control-to-manage-foreground-and-background-processes}{\texttt{How To Use Bash's Job Control to Manage Foreground and Background Processes}}
    \item \href{https://www.geeksforgeeks.org/linux-unix/chaining-commands-in-linux/}{\texttt{Chaining Commands in Linux}}
    \item \href{https://ubuntu.com/tutorials/command-line-for-beginners}{\texttt{The Linux command line for beginners}}
\end{itemize}

% Updated Footer
\vspace{\fill}
\begin{center}
\color{muted}
All rights reserved for \texttt{TogetherTeam} [NOT For Commercial Use] !!\\
contact the author via email \href{mailto:khaledmhfz2004@gmail.com}{\texttt{khaledmhfz2004@gmail.com}}
\end{center}

\end{document}